% \begin{lstlisting}[language=python]
%     print("hiu")
% \end{lstlisting}

% \lstdefinestyle{console}{
%     basicstyle=\ttfamily\small,
%     breaklines=true,
%     breakatwhitespace=false,
%     columns=fullflexible,
%     keepspaces=true,
%     postbreak=\mbox{\textcolor{gray}{$\hookrightarrow$}\space},
%     numbers=left,
%     numberstyle=\tiny,
%     frame=single
% }



% \lstset{language=Python, numbers=left, breaklines=true}
% \lstinputlisting{Listings/main.py}
% \inputminted[linenos=true, breaklines=true]{python}{Listings/main.py}
\codefromfile{main.py}{python}

\newpage
\textbf{Вывод программы:}

\VerbatimInput[
    fontsize=\small,
    breaklines=true,
    breakanywhere=true,
    frame=single
]{Listings/out.txt}

\textbf{Сравнение результатов:}

Для метода прямоугольников, трапеций и Симпсона были использованы порядки точности 2,2,4 соответсвенно.
При этом, для достижения погрешности $\epsilon = 10^{-3}$, потребовалось:
\begin{enumerate}
    \item $n=64$ разбиений для метода прямоугольников;
    \item $n=64$ для метода трапеций;
    \item $n=16$ для метода Симпсона.
\end{enumerate}